



In \exrange{exercise-4.1.1}{exercise-4.1.4}, rewrite a given expression as a power of $2$; that is, express it in the form $2^m$ for some $m$.

\begin{exercise}\label{exercise-4.1.1}
$\displaystyle \frac{4}{2^4}$ 
\end{exercise}

\begin{solution}
$2^{-2}$
\end{solution}

\begin{exercise}
$(2^3)^{-2}$ 
\end{exercise}

\begin{solution}
$2^{-6}$
\end{solution}

\begin{exercise}
$\displaystyle \frac{8^2}{2}$ 
\end{exercise}

\begin{solution}
$2^5$
\end{solution}

\begin{exercise}\label{exercise-4.1.4}
$\displaystyle \left(\frac{2^3}{2^4\cdot 2^5}\right)^2$
\end{exercise}

\begin{solution}
$2^{-12}$
\end{solution}







\noindent In \exrange{exercise-4.1.5}{exercise-4.1.8}, rewrite a given expression as a power of $5$; that is, express it in the form $5^m$ for some $m$.


\begin{exercise}\label{exercise-4.1.5}
$5^2\cdot 5^{-3}$  
\end{exercise}

\begin{solution}
$5^{-1}$
\end{solution}

\begin{exercise}
$\displaystyle \left(\frac{25^2}{5}\right)^{-1}$ 
\end{exercise}

\begin{solution}
$5^{-3}$
\end{solution}

\begin{exercise}
$\displaystyle \left(\frac{5^{-1}}{5^{-2}}\right)^2$  
\end{exercise}

\begin{solution}
$5^{2}$
\end{solution}

\begin{exercise}\label{exercise-4.1.8}
$\displaystyle \left(\frac{5^3}{25^2}\right)^{-2}$
\end{exercise}

\begin{solution}
$5^2$
\end{solution}







\noindent In \exrange{exercise-4.1.9}{exercise-4.1.15}, simplify a given expression if possible. If not possible, state so.


\begin{exercise}\label{exercise-4.1.9}
$(x^2y^5)^0$
\end{exercise}

\begin{solution}
$1$
\end{solution}

\begin{exercise}
$\displaystyle \left(\frac{x^{-3}}{xy^2}\right)^{-1}$
\end{exercise}

\begin{solution}
$x^4y^2$
\end{solution}

\begin{exercise}
$\displaystyle \frac{(x-y)^3}{x}$
\end{exercise}

\begin{solution}
Cannot be simplified.
\end{solution}

\begin{exercise}
$(2x^2y^4)^3$
\end{exercise}

\begin{solution}
$8x^6y^{12}$
\end{solution}

\begin{exercise}
$\displaystyle \frac{x+y}{y}$
\end{exercise}

\begin{solution}
Cannot be simplified.
\end{solution}

\begin{exercise}
$\displaystyle \left(\frac{\pi (x^{2})^6}{x^{2}y^{-4}}\right)^2$
\end{exercise}

\begin{solution}
$\pi^2x^{20}y^8$
\end{solution}

\begin{exercise}\label{exercise-4.1.15}
$\displaystyle \frac{x^2+y^3}{xy}$
\end{exercise}

\begin{solution}
Cannot be simplified.
\end{solution}





\begin{exercise}
Let $p$ and $r$ be positive integers. Use the definition of exponentiation
\[
a^m=\underbrace{a\cdot a\cdot  \ldots \cdot a}_{m \text{ times}}
\]
to explain why the following formula is valid:
\[
(a^p)^r=a^{p\cdot r}
.\]
\end{exercise}

\begin{solution}
\begin{align*}
(a^p)^r & = \left( \underbrace{a \cdot a \cdot \ldots \cdot a}_{p \text{ times}}\right)^r \\
& = \underbrace{\left( \underbrace{a \cdot a \cdot \ldots \cdot a}_{p \text{ times}}\right) \cdot \ldots \cdot \left( \underbrace{a \cdot a \cdot \ldots \cdot a}_{p \text{ times}}\right)}_{r \text{ times}} \\
&= \underbrace{a \cdot a \cdot \ldots \cdot a}_{p \cdot r \text{ times}} \\
&= a^{p\cdot r}
\end{align*}
\end{solution}
